\section[Cote 145 : l'exposant unité dans le groupe profini libre]{Cote 145 : l'exposant unité dans le groupe profini libre\protect\verifie{Folder145}}\label{sec:145}

La cote 145 étudie les automorphismes du groupe profini libre $\hat\pi$ sur $l_0$, $l_1$, avec $l_\infty l_1 l_0 = 1$, qui préservent les sous-groupes de lacets $L_0$, $L_1$, $L_\infty$ à des \q{rectificateurs} près. Les pages~34 et~35 caractérisent dans le cas discret, où l'exposant $p$ vaut $\pm 1$, les systèmes $(p, \sigma, 1, 1, 1)$ à rectificateurs triviaux : ce sont les six $(\operatorname{sg}\sigma, \sigma)$ du groupe $\Gamma_P$. La page conjecture (\q{Sans doute}) qu'il en va de même dans le cas profini, ramène la question à l'énoncé (14) ci-dessous et le laisse en \q{Exercice !} (\lot{145}{2}{34}). La lecture modernisée en donne une solution, qui est de l'édition, par un quotient diédral.

\noindent\emph{Conventions.} Mathlib n'a ni le groupe profini libre ni les puissances d'exposant dans $\hat{\mathbf Z}$. Une relation vaut dans $\hat\pi$ si et seulement si elle vaut pour toute image de $(l_0, l_1)$ dans tout groupe fini, et $p \in \hat{\mathbf Z}$ agit sur un groupe fini $G$ par son résidu modulo $\operatorname{card} G$. On représente donc $p$ par une famille compatible de résidus $p_n \in \mathbf Z/n$, $n \geq 1$, et l'on pose $x^p = x^{p_{\operatorname{card} G}}$ dans un groupe fini~$G$. La relation (14) pour $p$ signifie : pour tout groupe fini $G$ et tous $x, y, z \in G$ tels que $zyx = 1$, on a $z^p y^p x^p = 1$.

\begin{theoreme}[Grothendieck, pages 34 et 35, l'exercice\piste{145-unit-exponent-exercise}]\label{thm:145-exercice}
Soit $p \in \hat{\mathbf Z}^*$. Si $l_\infty^p\, l_1^p\, l_0^p = 1$ dans le groupe profini libre sur $l_0$, $l_1$, avec $l_\infty l_1 l_0 = 1$, alors $p = 1$. Autrement dit, la relation (14) vaut pour $p$ si et seulement si $p_n = 1$ pour tout $n \geq 1$.
\end{theoreme}

\begin{proof}
La réciproque est claire. Soit $N \geq 1$. Comme $p_2$ est inversible dans $\mathbf Z/2$, $p_2 = 1$, et par compatibilité le représentant $m \in [0, 2N[$ de $p_{2N}$ est impair. Dans le groupe diédral $D_N$, d'ordre $2N$, soient $x$ et $y$ deux réflexions dont le produit $r = yx$ est une rotation d'ordre $N$, et $z = r^{-1}$, de sorte que $zyx = 1$. Comme $m$ est impair, $x^m = x$ et $y^m = y$, donc $z^m y^m x^m = r^{-m} r = r^{1 - m}$, qui vaut $1$ si et seulement si $m \equiv 1 \pmod N$. La relation (14) dans $D_N$ donne donc $p_N = 1$, image de $p_{2N}$.
\end{proof}

\begin{remarque}
L'énoncé tient, et plus généralement : seule sert l'imparité de $p$, c'est-à-dire $p_2 = 1$, non l'inversibilité de $p$ aux nombres premiers impairs. L'imparité ne peut être omise, car $p = 0$ vérifie (14). La compatibilité de la famille ne sert qu'à passer de $p_{2N}$ à $p_N$ et à $p_2$.
\end{remarque}

Ne sont pas démontrés ici : la conséquence pour $S\hat{\hat{\mathcal E}}$ (les systèmes à rectificateurs triviaux sont les six $(\operatorname{sg}\sigma, \sigma)$ de $\Gamma_P$), qui demande le groupe $\Gamma_P$ et la réduction de la page aux $\sigma$ paires, ni le cas discret de la page~35.
